\documentclass[11pt]{article}
\usepackage[a4paper,margin=25mm]{geometry}
\usepackage{amsmath,amssymb,amsthm,hyperref}
\newtheorem{lemma}{Lemma}\newtheorem{theorem}{Theorem}
\newcommand{\E}{\mathbb E}
\title{Global spectral truncation for the direct marked identity}
\author{Complete analytic argument; independently reviewed}
\date{30 September 2026}
\begin{document}\maketitle
\noindent\textbf{Status.} Full independent mathematical reviews by the root and size-bounds agents passed. A separate independent-directions review also passed, including the appended unmarked and off-annulus results. This is not formal verification.

The important truncation is in \emph{differential order}, before the
infinite spectral tail is estimated. No estimate for the complete
high-order polarization of the infinite Poisson kernel is asserted.

Use the notation of
\path{independent_directions/marked_cancellation_identity.tex}:
$m=M+1$, $\ell=\lfloor M/4\rfloor$, $N=M-\ell$,
$z_j=\ell(1-p_j)$, and $s=\sum_{j\ne i}(1-p_j)$.
The row weights $w$ form a probability measure. A uniformly chosen
$\ell$-subset $B$ of the unmarked columns has weight
$A_B=\prod_{j\in B}p_j$, complement $R$, and mean
$v=N^{-1}\sum_{j\in R}z_j$. Put
\[
 a_k=\frac{e_k((z_j-v)_{j\in R})}{\binom Nk},\qquad
 a_0=1,\quad a_1=0.
\]
All expectations below include the uniformly chosen subset before
multiplication by $A_B$; thus no division by its possibly small mean
is present. We use the established finite-balance and row-mass bounds
in the following exact form, with constants independent of $m$:
\begin{equation}\label{inputs}
 \max_{j\in[m]}z_j\le C_0(s+1),\qquad
 \E_w e^{-b s}\le C_b/m\quad(b>0\text{ fixed}).
\end{equation}
These follow from the reviewed endpoint-balance and row-mass lemmas,
including after deleting one column. For the marked column itself,
the full-family estimate is rearranged to absorb its own contribution
to the deficit sum, for all sufficiently large $m$.

For a smooth function $F$ define the finite differential functional
\[
 \mathcal D_J F=(z_i-v)F(v)
 -\sum_{k=2}^{J}
 \left[\frac{(v-z_i)F^{(k)}(v)}{k!}
            +\frac{F^{(k-1)}(v)}{(k-1)!}\right]a_k.
\]
The exact marked identity says
$\E_w\E_B A_B\mathcal D_{d+1}P=0$ for every polynomial $P$ of
degree $d\le N-1$.

\begin{lemma}[Global integrated spectral derivative bound]\label{term}
Fix $L<\infty$. There is an absolute $C$ (depending on the constants
in \eqref{inputs}) such that, for every $j\ge0$ and $0\le k\le N$,
\begin{align}
 &(\ell+1)\E_w\E_B A_B|a_k|
 \left[(1+|z_i-v|)\frac{|L_j^{(k)}(v)|}{k!}
 +\boldsymbol1_{k\ge1}\frac{|L_j^{(k-1)}(v)|}{(k-1)!}\right]
 \notag\\[-2mm]
 &\hspace{30mm}\le C(k+1)^3
       \left(C\sqrt{\frac{j+1}{m}}\right)^k .             \label{termineq}
\end{align}
Derivatives above degree $j$ are zero.
\end{lemma}
\begin{proof}
Write $\sigma^2=\sum_{j\in R}(z_j-v)^2$. By \eqref{inputs},
\[
 \sigma^2\le\sum_{j\in R}z_j^2\le C(s+1)Nv.
\]
The global complex-variable coefficient estimate for centered
variables gives
\[
 |a_k|\le\left(\frac{e k\sigma^2}{N^2}\right)^{k/2}
 \le\left(\frac{Ck(s+1)v}{N}\right)^{k/2}.                \tag{3}
\]
For $v>0$ the reviewed Gaussian-displacement estimate is
\[
 |L_j^{(h)}(v)|\le e^{v/2}(j/v)^{h/2}\quad(0\le h\le j).
                                                                    \tag{4}
\]
This is also proved in
\path{size_bounds/gaussian_laguerre_variance_lemmas.tex}.
Using $j+1$ in place of $j$ includes all boundary cases.
If $v=0$, all $z_j$ in $R$ vanish, hence $a_k=0$ for $k\ge1$;
the estimates obtained by multiplying (3) and (4) extend to that case
by continuity.

For each subset $v\le\ell s/N$. Maclaurin's inequality gives
$\E_BA_B\le(1-s/M)^\ell\le e^{-\ell s/M}$. Consequently
\[
 \E_B A_B e^{v/2}\le e^{-c s}
\]
for an absolute $c>0$ and sufficiently large $m$, since
$\ell/M-\ell/(2N)$ stays positive. The Laplace bound in
\eqref{inputs} implies, for every real $h\ge0$,
\[
 (\ell+1)\E_w e^{-c s}(s+1)^h
       \le C^{h+1}\Gamma(h+1).                            \tag{5}
\]
For example, reserve half the exponential, bound the other half times
$(s+1)^h$ by $(C(h+1))^h$, and use Stirling's inequality.

For the first derivative term in \eqref{termineq}, (3)--(5), together
with $1+|z_i-v|\le C(s+1)$, give the factor
\[
 \left(C\sqrt{\frac{j+1}{m}}\right)^k
       \frac{k^{k/2}\Gamma(k/2+2)}{k!}
 \le C(k+1)^2\left(C\sqrt{\frac{j+1}{m}}\right)^k.
\]
For the second term, one uncancelled factor $v^{1/2}$ remains;
use $v\le C s$ and (5). Its bound is at most
\[
 C(k+1)^3m^{-1/2}
       \left(C\sqrt{\frac{j+1}{m}}\right)^{k-1},
\]
which is absorbed into \eqref{termineq}, since $j+1\ge1$.
The case $k=0$ follows directly from (4)--(5), and $k=1$ is zero.
\end{proof}

\begin{theorem}[Arbitrarily small global residual]\label{truncation}
Fix $0<\delta<L<\infty$ and an exponent $A>0$.
There are constants $\kappa>0$, $B>0$, and $C_A>0$ such that,
for all sufficiently large $m$, put
\[
 d=\lfloor\kappa m\rfloor,\qquad
 J=\lceil B\log m\rceil,\qquad
 \tau=C_A\log m/m.
\]
Then $J\le d<N$, and uniformly for $a\in[\delta,L]$,
\begin{equation}\label{conclusion}
 \left|(\ell+1)\E_w\E_B A_B\,
          \mathcal D_J K_\tau(\cdot,a)\right|\le m^{-A}.
\end{equation}
Here $K_\tau$ is the full positive Laguerre Poisson kernel. The
conclusion is global over all rows and all subsets, with no compact
restriction on $v$ and no exclusion of zero entries.
\end{theorem}
\begin{proof}
Let $r=1-\tau$ and
\[
 P_d(x)=\sum_{j=0}^d r^jL_j(a)L_j(x).
\]
The Hille--Hardy identity says
$K_\tau(x,a)=\sum_{j\ge0}r^jL_j(a)L_j(x)$.
The Gaussian estimate with $h=0$ gives $|L_j(a)|\le e^{L/2}$.
The marked identity applies exactly to $P_d$.

Choose $\kappa$ first, sufficiently small that the geometric ratio
in \eqref{termineq} for $j\le d$ is at most some $\rho<1$.
The contribution of differential orders $k>J$ in
$\mathcal D_{d+1}P_d$ is then bounded by
\[
 \frac C\tau\sum_{k>J}(k+1)^3\rho^k.
\]
The term $k=d+1$ is allowed because $d+1\le N$; derivatives of
$P_d$ above its degree simply vanish. Taking $B$ sufficiently large
makes this bound at most $\tfrac12m^{-A}$, even after allowing
$\tau^{-1}\le m$.

It remains to replace $P_d$ by $K_\tau$ only at orders $k\le J$.
For each such order, Lemma~\ref{term} and $r^j\le e^{-\tau j}$
bound the integrated absolute spectral tail by
\[
 C(k+1)^3\sum_{j>d}e^{-\tau j}
                   \left(C\sqrt{\frac{j+1}{m}}\right)^k.
\]
Splitting $e^{-\tau j}$ into two factors and summing by comparison
with the gamma integral gives, for $0<\tau<1/2$,
\begin{align*}
 \sum_{j>d}e^{-\tau j}(j+1)^{k/2}
 &\le e^{-\tau d/2}C^{k+1}
               \Gamma(k/2+1)\tau^{-k/2-1}.
\end{align*}
Consequently the preceding bound is at most
\[
 \frac{Ce^{-\tau d/2}}\tau (k+1)^4
       \left(C\sqrt{\frac{k+1}{m\tau}}\right)^k.          \tag{6}
\]
Choose $C_A$ sufficiently large compared with $B$ that the last ratio
is at most $1/2$ for all $k\le J$ and large $m$.
The sum of (6) is therefore at most $Ce^{-\tau d/2}/\tau$.
Increasing $C_A$ further so that $C_A\kappa$ exceeds a fixed multiple
of $A+2$ makes this at most $\tfrac12m^{-A}$.
All spectral sums of the finitely many derivatives used here are
absolutely integrable by the same estimates, justifying their
interchange with the two expectations.
Combining the high-order truncation, the exact identity for $P_d$,
and the low-order spectral replacement proves \eqref{conclusion}.
\end{proof}


\begin{theorem}[Unmarked companion]\label{unmarked}
With the choices in Theorem~\ref{truncation}, define
\[
 \mathcal D_J^0F=\sum_{k=0}^J\frac{F^{(k)}(v)}{k!}a_k.
\]
Let $\beta_\ell$ be the probability measure on $[0,\ell]$ with density
\[
 b_\ell(x)=\frac{\ell+1}{\ell}
             (1-x/\ell)^\ell\boldsymbol1_{0\le x\le\ell}.
\]
Then uniformly for $a\in[\delta,L]$,
\[
 \left|(\ell+1)\E_w\E_BA_B\mathcal D_J^0K_\tau(\cdot,a)
                       -\int K_\tau(x,a)\,d\beta_\ell(x)\right|
 \le m^{-A}.
\]
Changing the constants to obtain the displayed coefficient one is
understood. The densities $b_\ell$ are bounded above and below by
positive constants on every fixed compact interval for large $m$.
\end{theorem}
\begin{proof}
For any polynomial $P$ of degree at most $N$, the exact unmarked
success/failure identity gives
\[
 (\ell+1)\E_w\E_B A_B\operatorname{Pol}_N(P)(z_R)
       =(\ell+1)\int_0^1t^\ell P(\ell(1-t))\,dt
       =\int P\,d\beta_\ell.
\]
Apply it to $P_d$. Lemma~\ref{term} bounds its differential-order
truncation and the low-order spectral replacement on the row side,
exactly as in Theorem~\ref{truncation}.
On the reference side $b_\ell(x)\le2e^{-x}$ for large $m$, so
$\int e^{x/2}d\beta_\ell\le4$. Consequently the reference spectral
tail is at most $Ce^{-\tau d}/\tau$. Increasing the constants in the
previous proof, or running that proof with exponent $A+1$, completes
the estimate. The last assertion follows directly from the displayed
density and $(1-x/\ell)^\ell\to e^{-x}$ uniformly on compact sets.
\end{proof}

\begin{lemma}[Uniform complex bound]\label{complexglobal}
For $a\in[\delta,L]$, $0<\tau<1/2$, $v\ge0$, and $|\zeta|\le\tau$,
\[
 |K_\tau(v+\zeta,a)|\le C_L/\tau.
\]
Thus $|\partial_v^kK_\tau(v,a)|/k!\le C_L\tau^{-k-1}$.
\end{lemma}
\begin{proof}
The entire power series gives $|I_0(z)|\le e^{|z|}$.
Write $r=1-\tau$ and $y=v+\tau$. The exponent in this bound for
$K_\tau(v+\zeta,a)$ is at most
\[
 \frac{-ry-ra+2r\tau+2\sqrt{ray}}\tau
 \le\frac{a-ra+2r\tau}\tau\le L+2,
\]
where $-ry+2\sqrt{ray}\le a$ follows by completing a square.
Cauchy's formula proves the derivative assertion.
\end{proof}

\begin{theorem}[Off-annulus differential tails]\label{outside}
For every fixed $0<\delta<L$ there are constants $0<s_-<s_+<\infty$
such that, with the choices above, the contribution to either
\[
 (\ell+1)\E_w\E_B A_B|\mathcal D_JK_\tau|,
 \qquad
 (\ell+1)\E_w\E_B A_B|\mathcal D_J^0K_\tau|
\]
from rows $s\notin[s_-,s_+]$ is at most $m^{-A}$ for sufficiently
large $m$. The same statement holds if absolute values are placed on
every summand of either differential functional separately.
\end{theorem}
\begin{proof}
For $v\le\delta/16$ and $|\zeta|\le\delta/16$, the elementary bound
$|I_0(z)|\le e^{|z|}$ gives
\[
 |K_\tau(v+\zeta,a)|\le\tau^{-1}e^{-c\delta/\tau}.
\]
Indeed its exponent is at most
$[r\delta/16-ra+2\sqrt{ra\delta/8}]/\tau$, which is bounded above
by $-c\delta/\tau$ for $a\ge\delta$ and sufficiently small $\tau$.
For $v\ge16L$ and $|\zeta|\le v/16$, the same argument gives
\[
 |K_\tau(v+\zeta,a)|\le\tau^{-1}e^{-cv/\tau}.
\]
Cauchy's formula gives the corresponding derivative bounds, with
radii $\delta/16$ and $v/16$ respectively.

Take $s_-$ sufficiently small that $v\le\ell s/N\le\delta/16$
whenever $s\le s_-$. By \eqref{inputs} the row range and marked
value are bounded on those rows. The coefficient estimate
$|a_k|\le(CH\sqrt{k/N})^k$ therefore makes the sum of the Cauchy
bounds through $J=B\log m$ at most
$C\tau^{-1}e^{-c\delta/\tau}$. The marked factors are bounded as
well. Multiplication by $\ell+1\le m$ leaves a quantity smaller than
any prescribed inverse power of $m$.

For $s\ge s_+$, choose $s_+$ large enough that
$x=(\ell/M)s$ satisfies $x/2\ge16L$.
On subsets with $v\ge x/2$, the row range and marked value are
$O(s+1)=O(v)$ by \eqref{inputs}. The large-$v$ Cauchy radius thus
cancels the range in the coefficient estimate. Summing through $J$
gives at most a polynomial factor in $m$ times $e^{-cs/\tau}$,
again negligible even after integration against the probability $w$.

For completeness the exceptional subsets $v<x/2$ have uniform
probability at most $e^{-cm}$. To see this without assumptions on the
success weights, sample the $N$-subset $R$ uniformly. Maclaurin's
inequality applied to the positive numbers $e^{t(z_j-x)/N}$, followed
by the elementary bounded-variable exponential estimate, gives
\[
 \E_R e^{t(v-x)}\le e^{t^2H^2/(8N)}.
\]
Here $H$ is the unmarked population range. Hence
$\Pr(v<x/2)\le\exp[-Nx^2/(2H^2)]\le e^{-cm}$, because
$H\le C(s+1)$ and $s\ge s_+>0$.
On these subsets use Lemma~\ref{complexglobal}, not a full
high-order polarization bound. Since $s\le M$, all $z_j$ are
$O(m)$, and for $k\le J=B\log m$ the product of the crude coefficient
bound and $\tau^{-k-1}$ is at most $\exp(C_B(\log m)^2)$.
The marked prefactors and $\ell+1$ are polynomial in $m$, while
$A_B\le1$. Thus the exceptional-subset contribution is at most
$e^{-cm}\exp(C_B(\log m)^2)$, which is negligible. This proves all
assertions.
\end{proof}

\paragraph{Order of choices and scope.}
The constants are chosen in the order $\kappa$, $B$, then $C_A$.
They do not depend on $m$. The theorem supplies a global analytic
residual for the marked identity. Turning it into a sharp profile
bound still requires positivity/localization of the full kernel,
control of the success-weighted subset mean, and a monotonicity
argument. None of those subsequent conclusions is asserted here.
\end{document}
